\documentclass[11pt,a4paper]{article}
\usepackage[utf8]{inputenc}
\usepackage[T1]{fontenc}
\usepackage[french]{babel}
\usepackage{lmodern,geometry,microtype}
\usepackage{amsmath}
\usepackage{siunitx}
\usepackage[version=4]{mhchem}
\usepackage{xcolor,listings,hyperref}
\geometry{margin=2.4cm}
\sisetup{locale=FR,per-mode=symbol}
\DeclareSIUnit\molecule{molécule}
\DeclareSIUnit\atm{atm}
\lstset{language=Python,basicstyle=\ttfamily\small,backgroundcolor=\color{black!4},frame=single,breaklines=true,showstringspaces=false}
\hypersetup{colorlinks=true,linkcolor=blue!45!black,urlcolor=blue!55!black}

\title{Documentation de \texttt{2\_sections\_efficaces\_gaz\_hitran.py}}
\author{Projet climat 2026}
\date{\today}

\begin{document}
\maketitle
\tableofcontents

\section{Rôle du programme}

Ce fichier est le moteur spectral du modèle. Il calcule les sections efficaces
HITRAN de \ce{H2O}, \ce{CO2}, \ce{O3} et \ce{CH4} entre 4 et
\SI{100}{\micro\metre}. Il est à la fois exécutable pour afficher les quatre
spectres et importé par le calcul d'émissivité.

\section{Installation préalable}

\begin{lstlisting}[language=bash]
python3.12 -m venv .venv
source .venv/bin/activate
python -m pip install -r requirements.txt
\end{lstlisting}

Les paquets externes sont \texttt{numpy}, \texttt{matplotlib} et
\texttt{hitran-api}, importé sous le nom \texttt{hapi}. \texttt{functools} et
\texttt{pathlib} font partie de Python. Un accès réseau est nécessaire si les
tables HITRAN locales sont absentes.

\section{Imports du fichier}

\begin{lstlisting}
from functools import lru_cache
from pathlib import Path

import matplotlib.pyplot as plt
import numpy as np
from hapi import absorptionCoefficient_Voigt, db_begin, fetch
\end{lstlisting}

\texttt{lru\_cache} évite de répéter l'initialisation, les téléchargements et
certains calculs spectraux. Les autres imports assurent la gestion des chemins,
le calcul vectoriel, les tracés et l'accès à HITRAN.

\section{Constantes et conventions}

\begin{itemize}
  \item \texttt{DOSSIER\_DONNEES} pointe vers \path{donnees_hitran/} ;
  \item la grille couvre $\lambda\in[4,100]$~\si{\micro\metre} ;
  \item cela correspond à
        $\widetilde\nu\in[100,2500]$~\si{\per\centi\metre} ;
  \item les identifiants sont \texttt{H2O=1}, \texttt{CO2=2},
        \texttt{O3=3} et \texttt{CH4=6} ;
  \item seul l'isotopologue~1 de chaque gaz est demandé.
\end{itemize}

\section{Fonctions}

\subsection{Conversion spectrale}

\texttt{convertir\_lambda\_m\_en\_nombre\_onde\_cm\_1(longueur\_onde\_m)}
applique
\begin{equation}
  \widetilde\nu=\frac{1}{100\lambda}.
\end{equation}
Elle accepte un scalaire ou un tableau NumPy.

\subsection{Initialisation de HAPI}

\texttt{initialiser\_base\_hitran()} appelle \texttt{db\_begin} avec le cache
local. Le décorateur \texttt{@lru\_cache(maxsize=1)} limite l'initialisation à
un appel par processus.

\subsection{Téléchargement des raies}

\texttt{telecharger\_raies\_gaz(nom\_table, numero\_molecule,
nombre\_onde\_min, nombre\_onde\_max)} appelle \texttt{fetch} si le couple de
fichiers \path{.data}/\path{.header} est incomplet. Une table déjà présente
n'est pas contrôlée quant à sa couverture spectrale : si les bornes changent,
il faut vérifier ou recréer le cache.

\subsection{Calcul du profil de Voigt}

\texttt{charger\_section\_gaz(...)} reçoit le nom de table, l'identifiant
moléculaire, les bornes, le pas, la température et la pression en atmosphères.
Par défaut, le pas est \SI{0.05}{\per\centi\metre},
$T=\SI{296}{\kelvin}$ et $p=\SI{1}{\atm}$.

\textbf{Retour :} \texttt{(nombres\_onde, sections\_m2)}. Le facteur $10^{-4}$
convertit les \si{\centi\metre\squared\per\molecule} de HAPI en
\si{\metre\squared\per\molecule}. Le cache conserve huit combinaisons
d'arguments au maximum.

\subsection{Interpolation}

\texttt{interpoler\_section\_gaz(grille\_longueurs\_onde\_m, nombres\_onde,
sections\_gaz)} transforme la grille en nombres d'onde puis emploie
\texttt{numpy.interp}. Hors plage, la section vaut zéro.

\subsection{Interface principale}

\texttt{calculer\_sections\_gaz(nom\_table, numero\_molecule, grille,
temperature=296.0, pression\_atm=1.0)} enchaîne calcul et interpolation. Elle
renvoie un tableau de sections de même taille que la grille fournie. C'est la
fonction utilisée par le modèle d'émissivité.

\subsection{Diagnostic graphique}

\texttt{tracer\_sections\_des\_quatre\_gaz()} crée une grille de 20000 points,
calcule les quatre gaz à \SI{296}{\kelvin} et \SI{1}{\atm}, puis affiche une
figure $2\times2$.

\section{Exécution et import par un autre fichier}

Diagnostic autonome :
\begin{lstlisting}[language=bash]
python 2_sections_efficaces_gaz_hitran.py
\end{lstlisting}

Le nom du fichier commence par un chiffre et ne peut pas apparaître dans un
\texttt{import} Python classique. Le modèle l'importe ainsi :
\begin{lstlisting}
from importlib import import_module
module = import_module("2_sections_efficaces_gaz_hitran")
calculer_sections_gaz = module.calculer_sections_gaz
\end{lstlisting}

Le dossier du script doit donc se trouver dans \texttt{sys.path}, ce qui est le
cas lorsqu'on lance les programmes depuis ce dossier.

\section{Sorties et limites}

L'exécution autonome produit seulement une figure ; l'appel comme module renvoie
des tableaux NumPy. Le profil de Voigt n'inclut pas explicitement dans ce code
les nuages, les aérosols, le continuum ni les isotopologues minoritaires. Une
interpolation ponctuelle sur une grille grossière peut aussi masquer des raies
très étroites.

\end{document}
